%% Springer Nature LaTeX template (sn-jnl.cls)
%% Target: Signal, Image and Video Processing (Springer)
%%
%% Unduh sn-jnl.cls dan sn-basic.bst dari:
%%   https://www.springernature.com/gp/authors/campaigns/latex-author-support
%%
%% VERSI INI: seluruh tabel eksperimen sudah terisi angka nyata dari
%% baseline_mandiri.ipynb. TODO yang tersisa hanya empat:
%%   1. Flowers102 (perlu ekstraksi ulang; berkas lama tersimpan terpotong)
%%   2. Ablasi ukuran p (perlu citra asli, tidak bisa dari Excel 32x32)
%%   3. Appendix dataset enam citra (disarankan dihapus)
%%   4. Data administratif (pendanaan, kontribusi penulis, tautan kode)
%%
%% CATATAN KONSISTENSI: Tabel ablasi level intensitas dihitung di lingkungan
%% berbeda dari tabel lain, sehingga baris pertamanya membaca 92.76 sedangkan
%% Tabel 3 membaca 92.99. Jalankan ulang ablasi itu di mesin Anda sebelum
%% submit. Sel kodenya ada di komentar tepat di atas tabel tersebut.

\documentclass[pdflatex,sn-basic]{sn-jnl}

\usepackage{graphicx}
\usepackage{multirow}
\usepackage{amsmath,amssymb,amsfonts}
%% amsthm sudah dimuat oleh sn-jnl.cls
\usepackage{mathrsfs}
\usepackage[title]{appendix}
\usepackage{xcolor}
\usepackage{textcomp}
\usepackage{manyfoot}
\usepackage{booktabs}
\usepackage{algorithm}
\usepackage{algorithmicx}
\usepackage{algpseudocode}

%%%%%%%%%%%%%%%%%%%%%%%%%%%%%%%%%%%%%%%%%%%%%%%%%%%%%%%%%%%%%%%%%%%%%%%%
%% Deklarasi environment teorema.
%% sn-jnl.cls TIDAK mendefinisikan ini sendiri -- di berkas contoh Springer
%% blok berikut ada tetapi dikomentari, sehingga penulis harus mengaktifkannya.
%% Tanpa blok ini muncul galat "No counter 'theorem' defined".
%% Gaya thmstyleone/two/three disediakan oleh sn-jnl.cls.

\theoremstyle{thmstyleone}%
\newtheorem{theorem}{Theorem}%          penomoran berlanjut
\newtheorem{proposition}[theorem]{Proposition}%
\newtheorem{lemma}[theorem]{Lemma}%
\newtheorem{corollary}[theorem]{Corollary}%

\theoremstyle{thmstyletwo}%
\newtheorem{example}{Example}%
\newtheorem{remark}{Remark}%

\theoremstyle{thmstylethree}%
\newtheorem{definition}{Definition}%
\newtheorem{assumption}{Assumption}%

%% Hapus makro ini dan seluruh pemakaiannya sebelum submit.
\newcommand{\TODO}[1]{\textcolor{red}{\textbf{[LENGKAPI: #1]}}}

\raggedbottom

\begin{document}

\title[ECDF Colour Descriptors for Image Clustering]{Seven Numbers Suffice:
Low-Dimensional ECDF Colour Descriptors for Image Clustering}

%% Tanda bintang pada \author menandai corresponding author.
%% Urutan penulis tetap: Husty Serviana Husain penulis pertama,
%% Sapto Wahyu Indratno corresponding author.
\author[1,2]{\fnm{Husty Serviana} \sur{Husain}}\email{serviana@upi.edu}
\author*[1]{\fnm{Sapto Wahyu} \sur{Indratno}}\email{sapto@itb.ac.id}
\author[1]{\fnm{Sandy} \sur{Vantika}}\email{sandy.vantika@itb.ac.id}
\author[3]{\fnm{Nugraha Priya} \sur{Utama}}\email{utama@staff.stei.itb.ac.id}

\affil[1]{\orgdiv{Department of Mathematics, Faculty of Mathematics and
Natural Sciences}, \orgname{Institut Teknologi Bandung},
\orgaddress{\city{Bandung}, \country{Indonesia}}}
\affil[2]{\orgdiv{Department of Mathematics Education, Faculty of Mathematics
and Natural Sciences Education}, \orgname{Universitas Pendidikan Indonesia},
\orgaddress{\city{Bandung}, \country{Indonesia}}}
\affil[3]{\orgdiv{School of Electrical Engineering and Informatics},
\orgname{Institut Teknologi Bandung},
\orgaddress{\city{Bandung}, \country{Indonesia}}}

\abstract{We ask how far colour information can be compressed before image
clustering degrades. Each image is represented by the values of its
per-channel empirical distribution function at three intensity levels, giving
a seven-dimensional descriptor derived from $32\times32\times3$ pixel values,
a compression of $439{:}1$. On 4400 rose images in four colour classes, a
linear support vector machine and linear discriminant analysis both reach
$100\%$ five-fold cross-validated accuracy on this descriptor, and $K$-means
clustering reaches $92.99\pm7.34\%$ over 100 random initialisations. The
descriptor outperforms colour moments ($75.21\pm12.90\%$), a quantised joint
RGB histogram of 64 dimensions ($80.02\pm13.20\%$), and features from an
ImageNet-pretrained ResNet-50 ($48.52\pm12.75\%$), the last by a wide margin
that we trace to a mismatch of inductive bias rather than to any general
superiority over learned representations. We prove that the descriptor is
exactly invariant under strictly increasing photometric transformations and
verify this numerically under gamma correction. We replace the common but
false claim that distinct colour distributions imply linear separability with
two valid statements: exact separability under a margin condition, and a
distribution-free Chebyshev bound whose ordering across class pairs matches
the observed confusion matrix. A rank transform of the descriptor leaves mean
accuracy unchanged ($p=0.41$, Cohen's $d_z=0.08$) but reduces the standard
deviation across initialisations from $7.34$ to $2.73$ percentage points, so
we recommend it for stability rather than accuracy. Finally, on a five-category colour subset of
Flowers102 every method fails, no descriptor exceeding $45.23\%$ and the
supervised ceiling reaching only $59.16\%$; there the same Chebyshev bound is
vacuous for all ten class pairs and the between-class to within-class colour
variance ratio falls from $7.481$ to $0.361$. We propose both quantities as
feasibility criteria computable before clustering.}

\keywords{Image clustering, Empirical cumulative distribution function, Colour
descriptor, Photometric invariance, Feature compression, $K$-means}

\maketitle

%%%%%%%%%%%%%%%%%%%%%%%%%%%%%%%%%%%%%%%%%%%%%%%%%%%%%%%%%%%%%%%%%%%%%%%%
\section{Introduction}\label{sec:intro}

Colour is the cheapest signal available to an image clustering system and often
the most informative one. Learned representations dominate current practice,
but they are opaque, costly to compute, and presuppose training data that many
applied settings lack. Statistical colour descriptors remain attractive
wherever interpretability, low latency, or auditability matter
\cite{ref:jain2010,ref:rudin2019}.

Our question is quantitative rather than architectural: how few numbers per
image suffice? Colour histograms with $b$ bins per channel yield $3b$ features
and depend on bin placement \cite{ref:swain1991}. Colour moments compress
further but discard the shape of the distribution \cite{ref:stricker1995}. The
empirical cumulative distribution function (ECDF) of a channel encodes the
entire colour distribution of an image, and evaluating it at a handful of
intensity levels gives a descriptor whose dimension the designer sets directly
\cite{ref:van2016}.

Three gaps recur in prior work that uses ECDF-based representations for image
analysis. First, the sampling model is usually left implicit, and several
published treatments invoke the Glivenko--Cantelli theorem for image pixels,
which are spatially correlated and hence not independent. Second, separability
is commonly asserted in the form ``distinct distributions imply linearly
separable features'', which is false in general. Third, transformations applied
to descriptors before clustering, typically rank-based or quantile
normalisations, are reported as improvements without controlling for the
initialisation variance of $K$-means.

This paper addresses all three, and reports two results we did not expect.

\begin{enumerate}
\item \textbf{Methodological.} A seven-dimensional descriptor built from
per-channel ECDFs at three intensity levels supports perfect linear class
separation on a 4400-image rose dataset under five-fold cross-validation. It
outperforms colour moments, a 64-dimensional RGB histogram, and
ImageNet-pretrained ResNet-50 features by margins reported in
Section~\ref{sec:baseline-results}.

\item \textbf{Theoretical.} We place the descriptor in a sampling model whose
randomness resides at the image level rather than the pixel level, so no
independence assumption on pixels is needed
(Section~\ref{sec:measures}). We prove exact invariance under strictly
increasing photometric transformations
(Proposition~\ref{prop:invariance}), and we replace the false separability
implication with an exact result under a margin condition together with a
distribution-free bound estimable from data
(Theorem~\ref{thm:separability}).

\item \textbf{Practical.} The between-class to within-class colour variance
ratio, computable in one pass and before any clustering, predicts whether a
colour-only descriptor can succeed. We also report two negative results: the
rank transform buys stability but not accuracy, and $K$-means is not the best
clustering algorithm for this descriptor.
\end{enumerate}

Section~\ref{sec:theory} develops the sampling model, the descriptor, the
invariance property, and the separability analysis.
Section~\ref{sec:method} describes the pipeline, the baselines, and the
evaluation protocol. Section~\ref{sec:experiments} reports the experiments.
Section~\ref{sec:discussion} discusses the failure regime and the feasibility
criterion, and Section~\ref{sec:conclusion} concludes.

%%%%%%%%%%%%%%%%%%%%%%%%%%%%%%%%%%%%%%%%%%%%%%%%%%%%%%%%%%%%%%%%%%%%%%%%
\section{Theoretical foundations}\label{sec:theory}

\subsection{Setup and notation}\label{sec:notation}

Let $\mathcal{I}=\{0,1,\dots,255\}$ be the intensity alphabet and
$\mathcal{C}=\{R,G,B\}$ the colour channels. A digital image is a map
$J:\{1,\dots,p\}^2\to\mathcal{I}^3$; we write $C^{(g)}_{ij}$ for the intensity
of channel $C$ at pixel $(i,j)$ in image $g$. Throughout, $N$ is the number of
images, $m$ the number of sampled intensity levels $t_1<\dots<t_m$, and $d$ the
descriptor dimension. In the experiments $p=32$, $m=3$ and
$(t_1,t_2,t_3)=(100,150,200)$.

\subsection{Images as empirical colour measures}\label{sec:measures}

\begin{definition}[Empirical colour measure]\label{def:measure}
For image $g$ and channel $C$, the empirical colour measure is the probability
measure on $\mathcal{I}$
\begin{equation}
\mu^{(g)}_C \;=\; \frac{1}{p^2}\sum_{i=1}^{p}\sum_{j=1}^{p}
\delta_{C^{(g)}_{ij}},
\end{equation}
where $\delta_a$ is the Dirac measure at $a$. Its distribution function is
\begin{equation}\label{eq:ecdf-image}
F^{(g)}_C(x) \;=\; \mu^{(g)}_C\big([0,x]\big)
\;=\;\frac{1}{p^2}\sum_{i=1}^{p}\sum_{j=1}^{p}\mathbf{1}\{C^{(g)}_{ij}\le x\},
\qquad x\in\mathcal{I}.
\end{equation}
\end{definition}

\begin{remark}[Statistical status of \eqref{eq:ecdf-image}]\label{rem:no-iid}
$F^{(g)}_C$ is not an estimator. It is the exact distribution function of the
finite population of $p^2$ pixels constituting image $g$, computed rather than
inferred. No independence assumption on the pixels is required, and none is
made. Pixels of a natural image are spatially correlated, so an independence
hypothesis at the pixel level would be untenable and the Glivenko--Cantelli
theorem cannot legitimately be invoked there. Randomness enters one level
higher, through the sampling of images, as formalised in
Assumption~\ref{ass:sampling}.
\end{remark}

\begin{assumption}[Sampling model]\label{ass:sampling}
Let $\mathcal{P}(\mathcal{I})$ be the set of probability measures on
$\mathcal{I}$. Each class $c\in\{1,\dots,K\}$ is identified with a probability
distribution $\Pi_c$ on $\mathcal{P}(\mathcal{I})^3$, the law of the triple
$(\mu_R,\mu_G,\mu_B)$ of a randomly drawn image of that class. Images of a
given class are independent draws from $\Pi_c$.
\end{assumption}

\begin{definition}[Descriptor]\label{def:feature}
The feature map $\Phi:\mathcal{P}(\mathcal{I})^3\to[0,1]^{d}$ evaluates channel
distribution functions at a set $S$ of channel--level pairs,
\begin{equation}
\Phi(\mu_R,\mu_G,\mu_B)=\big(F_C(t)\big)_{(C,t)\in S},\qquad d=|S|.
\end{equation}
We write $\mathbf{v}_g=\Phi(\mu^{(g)}_R,\mu^{(g)}_G,\mu^{(g)}_B)\in[0,1]^{d}$
with coordinates $v_{g,1},\dots,v_{g,d}$.
\end{definition}

Every coordinate of $\Phi$ is a value of a distribution function, so the
feature space is the unit cube $[0,1]^{d}$ and not $\mathbb{R}^{d}$.

\subsection{Photometric invariance}\label{sec:invariance}

\begin{proposition}[Invariance under monotone photometric
transformations]\label{prop:invariance}
Let $\varphi:\mathcal{I}\to\mathcal{I}$ be strictly increasing and let
$\widetilde{J}^{(g)}$ be obtained by applying $\varphi$ to every pixel of
$J^{(g)}$ in every channel, with distribution functions
$\widetilde{F}^{(g)}_C$. Then for every channel $C$ and level $t$,
\begin{equation}
\widetilde{F}^{(g)}_C\big(\varphi(t)\big)=F^{(g)}_C(t).
\end{equation}
Consequently, if the sampling grid is transformed accordingly,
$t\mapsto\varphi(t)$, the descriptors $\mathbf{v}_g$ and every clustering
computed from them are unchanged.
\end{proposition}

\begin{proof}
Because $\varphi$ is strictly increasing, $\varphi(a)\le\varphi(t)$ if and only
if $a\le t$. Hence
\begin{equation*}
\widetilde{F}^{(g)}_C(\varphi(t))
=\frac{1}{p^2}\sum_{i,j}\mathbf{1}\{\varphi(C^{(g)}_{ij})\le\varphi(t)\}
=\frac{1}{p^2}\sum_{i,j}\mathbf{1}\{C^{(g)}_{ij}\le t\}
=F^{(g)}_C(t).
\end{equation*}
The remaining claims follow because all subsequent computations depend on the
images only through $\mathbf{v}_1,\dots,\mathbf{v}_N$.
\end{proof}

Gamma correction $\varphi(a)=255\,(a/255)^{\gamma}$, monotone tone mapping and
per-channel monotone illumination changes are all of this form. Colour moments
and raw histogram bin counts are invariant under none of them.
Section~\ref{sec:invariance-exp} verifies the proposition numerically.

\subsection{Separability}\label{sec:separability}

\begin{remark}[Distinct distributions do not imply
separability]\label{rem:counterexample}
Two obstructions must be distinguished.

\emph{(a) Loss of information at the sampling grid.} Let $\mu_1\neq\mu_2$ be
colour measures with $\mu_1([0,t_k])=\mu_2([0,t_k])$ for $k=1,\dots,m$. Such
pairs exist whenever $m<255$, since $\Phi$ maps a $255$-dimensional simplex
into $[0,1]^{m}$ per channel. The classes then have identical descriptors and
no separating hyperplane exists.

\emph{(b) Within-class overlap.} Even when $\Phi$ separates the class means,
the laws of $\mathbf{v}$ under $\Pi_1$ and $\Pi_2$ may overlap in $[0,1]^{d}$,
so no hyperplane classifies every image correctly.
\end{remark}

\begin{theorem}[Separability]\label{thm:separability}
Let $\Pi_1,\Pi_2$ be two classes and $\mathbf{v}$ the descriptor of
Definition~\ref{def:feature}.
\begin{enumerate}
\item[(i)] \emph{(Exact separability under a margin condition.)} Suppose there
exist a coordinate $k$, a threshold $a\in(0,1)$ and $\gamma>0$ with
$\Pi_1(v_{k}\le a-\gamma)=1$ and $\Pi_2(v_{k}\ge a+\gamma)=1$. Then the
hyperplane $\{\mathbf{x}\in[0,1]^{d}:x_{k}=a\}$ separates the classes with
geometric margin at least $\gamma$.

\item[(ii)] \emph{(Approximate separability.)} Let
$\mu_{1,k}=\mathbb{E}_{\Pi_1}[v_k]$, $\mu_{2,k}=\mathbb{E}_{\Pi_2}[v_k]$,
$\Delta_k=|\mu_{1,k}-\mu_{2,k}|>0$ and
$\sigma_k^2=\max\{\mathrm{Var}_{\Pi_1}(v_k),\mathrm{Var}_{\Pi_2}(v_k)\}$. Then
the threshold rule at $a=(\mu_{1,k}+\mu_{2,k})/2$ misclassifies, under either
class, a fraction at most
\begin{equation}\label{eq:cheb}
\varepsilon_k \;=\; \min\Big\{1,\ \frac{4\sigma_k^{2}}{\Delta_k^{2}}\Big\}.
\end{equation}
\end{enumerate}
\end{theorem}

\begin{proof}
(i) Take $\mathbf{w}=e_k$ and $b=-a$. For $\mathbf{v}\sim\Pi_1$,
$\mathbf{w}^{\top}\mathbf{v}+b=v_k-a\le-\gamma$ almost surely; for
$\mathbf{v}\sim\Pi_2$, $\mathbf{w}^{\top}\mathbf{v}+b\ge\gamma$ almost surely.
Since $\|\mathbf{w}\|=1$ the geometric margin is at least $\gamma$.

(ii) Without loss of generality $\mu_{1,k}<\mu_{2,k}$. An image of class $1$ is
misclassified only if $v_k\ge a$, which forces $|v_k-\mu_{1,k}|\ge\Delta_k/2$.
By Chebyshev's inequality,
\begin{equation*}
\Pi_1(v_k\ge a)\le
\Pi_1\Big(|v_k-\mu_{1,k}|\ge\tfrac{\Delta_k}{2}\Big)
\le\frac{\mathrm{Var}_{\Pi_1}(v_k)}{(\Delta_k/2)^{2}}
\le\frac{4\sigma_k^{2}}{\Delta_k^{2}} .
\end{equation*}
The argument for class $2$ is symmetric, and probabilities never exceed one.
\end{proof}

\begin{remark}\label{rem:bound-use}
Statement (ii) is distribution-free and needs only finite second moments. Under
a boundedness or sub-Gaussian assumption the polynomial bound may be replaced
by an exponential one, at the cost of a stronger hypothesis. Since $\Delta_k$
and $\sigma_k$ are estimable, $\varepsilon_k$ can be computed for every class
pair and coordinate and used to predict which pairs will be confused;
Section~\ref{sec:bounds-exp} compares the induced ordering with the observed
confusion matrix.
\end{remark}

\subsection{The rank transform}\label{sec:rank}

\begin{definition}[Empirical marginal transform]\label{def:rank}
For coordinate $k$ and a sample of $N$ images, define
\begin{equation}\label{eq:rank}
\widehat{G}_{k}(s)=\frac{1}{N}\sum_{g=1}^{N}\mathbf{1}\{v_{g,k}\le s\},
\qquad s\in[0,1],
\end{equation}
and set $\widetilde{v}_{g,k}=\widehat{G}_{k}(v_{g,k})$.
\end{definition}

\begin{remark}[Terminology and relation to existing
transformations]\label{rem:pit}
$\widehat{G}_k$ is an empirical distribution function of the $N$ observed
values $v_{1,k},\dots,v_{N,k}$; it is not a cumulative distribution function of
a postulated random variable, and we do not describe it as one. The map
$v_{g,k}\mapsto\widehat{G}_k(v_{g,k})$ returns the normalised rank of image $g$
along coordinate $k$. It coincides with the empirical probability integral
transform \cite{ref:rosenblatt1952}, with rank-based normalisation, and, up to
the linear interpolation used between sample quantiles, with the
\texttt{QuantileTransformer} of scikit-learn \cite{ref:pedregosa2011}: across
our 4400-image dataset the average correlation between the two is $0.9908$. We
claim no novelty for the transformation and evaluate it as a baseline.
\end{remark}

\begin{proposition}[Elementary properties]\label{prop:properties}
Fix $k$ and suppose $v_{1,k},\dots,v_{N,k}$ are pairwise distinct. Then
$\widehat{G}_k$ is nondecreasing and right-continuous with
$\widetilde{v}_{g,k}\in\{1/N,\dots,1\}$; the map
$v_{g,k}\mapsto\widetilde{v}_{g,k}$ is nondecreasing; and the empirical
distribution of $\{\widetilde{v}_{g,k}\}_{g=1}^{N}$ is uniform on
$\{1/N,\dots,1\}$, with mean $(N+1)/(2N)$ and variance $(N^2-1)/(12N^2)$.
\end{proposition}

\begin{proof}
The first two claims follow from \eqref{eq:rank}. For the third, distinctness
makes $g\mapsto\widetilde{v}_{g,k}$ a bijection onto $\{1/N,\dots,1\}$; the
moments are those of a discrete uniform law.
\end{proof}

After the transform every coordinate carries the same empirical variance, so no
coordinate dominates the Euclidean distance used by $K$-means. This is a
statement about geometry, not accuracy. Section~\ref{sec:rank-eval} shows that
the practical consequence is reduced sensitivity to initialisation rather than
improved mean accuracy.

$\widehat{G}_k$ is a nondecreasing step function whose inverse is not unique.
No step of our argument requires a generalised inverse, so the question does
not arise.

\begin{lemma}[Glivenko--Cantelli at the image level]\label{lemma:gc}
Fix a coordinate $k$ and a class $c$, and let $G_{k,c}$ be the distribution
function of $v_{\cdot,k}$ under $\Pi_c$. Under Assumption~\ref{ass:sampling},
$\sup_{s\in[0,1]}|\widehat{G}_{k}(s)-G_{k,c}(s)|\to0$ almost surely as
$N\to\infty$.
\end{lemma}

\begin{proof}
By Assumption~\ref{ass:sampling} the values $v_{1,k},\dots,v_{N,k}$ are
independent and identically distributed with distribution function $G_{k,c}$,
and $\widehat{G}_k$ is their empirical distribution function. The classical
Glivenko--Cantelli theorem applies verbatim \cite{ref:wasserman2006}.
\end{proof}

The asymptotics are in the number of images $N$, which the experimental design
controls, and not in the number of pixels $p^2$, which is fixed by
preprocessing and for which the required independence fails.

\subsection{Computational complexity}\label{sec:complexity}

\begin{proposition}[Cost of the pipeline]\label{prop:complexity}
With $N$ images, $p^2$ pixels per image, descriptor dimension $d$, $K$ clusters
and $I$ iterations, the pipeline runs in time
$O(Np^{2}+dN\log N+dNKI)$ and uses $O(Nd)$ working memory beyond the input.
\end{proposition}

\begin{proof}
Because $|\mathcal{I}|=256$, the counts $\#\{(i,j):C^{(g)}_{ij}\le t\}$ for all
$t$ follow from a $256$-bin histogram and a prefix sum, costing $O(p^{2}+256)$
per channel and image, hence $O(Np^{2})$ overall. No sorting of pixels is
required, so no logarithmic factor in $p^2$ arises. Evaluating \eqref{eq:rank}
at all $N$ observed values of a coordinate requires one sort of $N$ numbers,
repeated over $d$ coordinates. Lloyd's algorithm costs $O(NdKI)$. Memory is
dominated by the $N\times d$ feature matrix.
\end{proof}

%%%%%%%%%%%%%%%%%%%%%%%%%%%%%%%%%%%%%%%%%%%%%%%%%%%%%%%%%%%%%%%%%%%%%%%%
\section{Method}\label{sec:method}

\subsection{Pipeline}\label{sec:pipeline}

\paragraph{Preprocessing.} Images are read in RGB and resized to $32\times32$
pixels by bilinear interpolation (\texttt{cv2.resize} with its default
\texttt{INTER\_LINEAR}). Resizing standardises the descriptor across
resolutions; because the descriptor depends on an image only through its colour
distribution, downsampling preserves the quantity of interest so long as it
introduces no systematic colour shift. Section~\ref{sec:ablation} confirms this
empirically over resolutions from $16\times16$ to $128\times128$.

\paragraph{Descriptor.} We evaluate the per-channel ECDF at intensity levels
$100$, $150$ and $200$. Two coordinates, $G_{100}$ and $B_{100}$, are
near-degenerate on our data and are dropped, giving the seven-dimensional
descriptor
\begin{equation}\label{eq:seven}
\mathbf{v}_g=\big(R_{100},R_{150},R_{200},G_{150},G_{200},B_{150},B_{200}\big),
\end{equation}
where $C_{t}$ abbreviates $F^{(g)}_C(t)$. Section~\ref{sec:ablation} compares
this choice against the full nine-coordinate grid and alternative level sets.
The levels were selected empirically; we make no claim of optimality or of
transfer to other datasets.

\paragraph{Clustering.} We run Lloyd's algorithm with $k$-means++
initialisation \cite{ref:arthur2007} and $K$ equal to the number of
ground-truth classes. No label information enters the initialisation. Cluster
identities carry no meaning, so before computing any label-dependent metric we
solve the assignment problem between clusters and classes with the Hungarian
algorithm on the contingency matrix \cite{ref:kuhn1955}.
Section~\ref{sec:algos} additionally reports Gaussian mixtures, agglomerative
clustering and DBSCAN on the same descriptor.

\subsection{Baselines}\label{sec:baselines}

We compare the descriptor against colour moments (mean, standard deviation and
skewness per channel, $d=9$) \cite{ref:stricker1995}, a quantised joint RGB
histogram with four bins per channel ($d=64$) \cite{ref:swain1991}, and
features from the penultimate layer of an ImageNet-pretrained ResNet-50
($d=2048$) \cite{ref:he2016resnet}. ResNet-50 features are extracted twice:
from the original-resolution images, which is the setting favourable to the
network, and from the $32\times32$ images upsampled to $224\times224$, which
equalises the input available to all methods. Descriptors with $d>20$ are
standardised before clustering.

We compare the untransformed descriptor against min--max scaling, $z$-score
standardisation, the rank transform of Definition~\ref{def:rank}, and
scikit-learn's \texttt{QuantileTransformer}.

\subsection{Evaluation protocol}\label{sec:protocol}

$K$-means is sensitive to initialisation, so single-run figures are
uninformative. We report two protocols.

\emph{Protocol A (stability).} The full dataset is clustered 100 times with
different random seeds. We report mean, standard deviation, minimum, maximum,
and the empirical $2.5$ and $97.5$ percentiles. These are percentile intervals
of the observed distribution across initialisations and require no
distributional assumption or pivotal quantity.

\emph{Protocol B (leakage-free generalisation).} Stratified five-fold
cross-validation repeated over 20 seeds, giving 100 runs. Any transformation is
fitted on the training folds only and applied to the held-out fold; the
clustering is fitted on the training folds and the held-out fold assigned by
nearest centroid. This is necessary because the rank transform is computed
across images and would otherwise leak information from the test fold.

We report external metrics (accuracy, precision, recall and $F_1$ after
Hungarian matching, adjusted Rand index \cite{ref:hubert1985}, normalised
mutual information) and internal metrics (silhouette \cite{ref:rousseeuw1987},
Davies--Bouldin \cite{ref:davies1979}, Calinski--Harabasz). Comparisons between
representations use paired $t$-tests across seeds, reported with sample size,
test statistic, $p$-value and Cohen's $d_z$.

Hyperparameters of all clustering baselines are selected by silhouette
coefficient, which uses no label information. Selecting them by accuracy would
leak labels into the baselines and invalidate the comparison.

%%%%%%%%%%%%%%%%%%%%%%%%%%%%%%%%%%%%%%%%%%%%%%%%%%%%%%%%%%%%%%%%%%%%%%%%
\section{Experiments}\label{sec:experiments}

\subsection{Datasets}\label{sec:datasets}

\emph{Rose-4400.} 4400 rose images in four colour classes, balanced at 1100
images per class: pink, red, white and yellow. The images are drawn from a
publicly available Kaggle collection \cite{ref:rosekaggle}. We verified that
our results reproduce from that source: re-downloading the collection and
re-running the extraction pipeline yields, for each of the four classes, a set
of descriptors identical to ours to machine precision, the two extractions
differing only in file naming. \TODO{cantumkan tanggal unduhan verifikasi}

\emph{Flowers102-C.} A subset of the Oxford Flowers102 dataset
\cite{ref:nilsback2008} of 181 images in five colour categories: red, blue,
yellow and purple with 40 images each, and brown with 21.

The colour labels are not human annotations, and we state their provenance in
full because it bears on how the results should be read. They were generated
automatically: the dominant colour of each image was obtained by $K$-means with
three clusters on a $100\times100$ version of the image, then mapped to one of
five categories by thresholds on hue and saturation, with a fallback rule on
RGB ratios where no threshold applied. Two properties of that rule matter. The
hue ranges assigned to blue ($100$--$140$) and purple ($130$--$170$) overlap and
are evaluated in order, so every image with hue between $130$ and $140$ is
labelled blue; pure purple, RGB $(128,0,255)$, is labelled blue by this rule.
The brown category requires hue in $[10,20]$, saturation in $[40,80]$ and value
below $180$ simultaneously, a conjunction rarely satisfied, which is why it
yields 21 images against a quota of 40 even when all 8189 images are scanned.
The category boundaries are therefore artefacts of the thresholds rather than
perceptual divisions.

This construction makes the labels a function of colour, which would inflate
any success a colour descriptor achieved here. It does not inflate a failure.
As Section~\ref{sec:failure} shows, no colour descriptor recovers these
categories, and the reason is a mismatch between two different measurements of
colour: the labels encode the dominant colour of one local region, while the
descriptor encodes the distribution over the whole image. The dataset therefore
serves as a negative control, not as an independent benchmark, and we do not
report it as evidence of generalisation.

\subsection{Linear separability of the descriptor}\label{sec:linear}

Table~\ref{tab:supervised} reports supervised five-fold cross-validated
accuracy on the descriptor of \eqref{eq:seven}. Both linear classifiers
separate the four rose classes perfectly. Since the descriptor derives from
$32\times32\times3=3072$ pixel values, this is a compression of $439{:}1$ with
no loss of linear discriminability.

\begin{table}[h]
\centering
\caption{Supervised five-fold cross-validated accuracy on the
seven-dimensional descriptor, Rose-4400.}\label{tab:supervised}
\begin{tabular}{@{}lc@{}}
\toprule
Classifier & Accuracy \\
\midrule
Linear SVM & $100.0\%$ \\
Linear discriminant analysis & $100.0\%$ \\
\botrule
\end{tabular}
\end{table}

\subsection{Comparison with baseline descriptors}\label{sec:baseline-results}

Table~\ref{tab:descriptors} reports Protocol A for all descriptors. The
seven-dimensional ECDF descriptor outperforms colour moments by $17.8$
percentage points and the 64-dimensional RGB histogram by $13.0$, while using
fewer dimensions than either.

\begin{table}[h]
\centering
\caption{Protocol A on Rose-4400: $K$-means with $k$-means++ over 100 seeds.
Accuracy after Hungarian matching. ARI is the adjusted Rand index, Silh.\ the
silhouette coefficient, DB the Davies--Bouldin index.}\label{tab:descriptors}
\begin{tabular}{@{}lcccccc@{}}
\toprule
Descriptor & $d$ & Accuracy & SD & ARI & Silh. & DB \\
\midrule
ECDF, seven coordinates       & $7$    & $\mathbf{92.99}$ & $7.34$  & $\mathbf{0.864}$ & $0.709$ & $0.528$ \\
ECDF, full grid               & $9$    & $89.85$ & $11.26$ & $0.840$ & $0.706$ & $0.492$ \\
RGB histogram, $4^3$ bins     & $64$   & $80.02$ & $13.20$ & $0.705$ & $0.507$ & $1.079$ \\
Colour moments                & $9$    & $75.21$ & $12.90$ & $0.655$ & $0.712$ & $0.444$ \\
ResNet-50, $32\times32$ input & $2048$ & $60.92$ & $11.52$ & $0.400$ & $0.327$ & $1.383$ \\
ResNet-50, original images    & $2048$ & $48.52$ & $12.75$ & $0.259$ & $0.370$ & $1.308$ \\
\botrule
\end{tabular}
\end{table}

The ResNet-50 result requires careful interpretation, and we state plainly what
it does and does not show. It does not show that a seven-number statistic beats
learned representations in general. It shows that a frozen
ImageNet-pretrained backbone carries the wrong inductive bias for this
particular task. The network was trained to recognise object categories, and
standard ImageNet augmentation includes colour jitter, which explicitly
encourages invariance to the very signal that defines our labels. Every image
in Rose-4400 is a rose; the clusters that ResNet-50 features induce are
organised by pose, background and texture, not by petal colour.

The two ResNet rows support this reading directly. Reducing the input to
$32\times32$ before upsampling destroys texture and fine spatial structure
while leaving the colour distribution intact, and it \emph{improves} accuracy
from $48.52\%$ to $60.92\%$, with the adjusted Rand index rising from $0.259$
to $0.400$. A baseline degraded by an information bottleneck would move the
other way. What the bottleneck removes is precisely the information that was
leading the network away from the labels.

We would expect a network fine-tuned on colour-defined labels, or a
self-supervised model trained without colour augmentation, to behave
differently. We did not run such experiments and make no claim about them.

\subsection{Comparison of representations}\label{sec:rank-eval}

Table~\ref{tab:representations} reports Protocol A across transformations of
the descriptor, Table~\ref{tab:protocolB} reports Protocol B, and
Table~\ref{tab:ttest} gives paired $t$-tests across the 100 seeds.

\begin{table}[h]
\centering
\caption{Protocol A: transformations of the seven-dimensional descriptor,
Rose-4400, 100 seeds. CI is the $2.5$--$97.5$ percentile interval across
initialisations.}\label{tab:representations}
\begin{tabular}{@{}lcccccc@{}}
\toprule
Representation & Accuracy & SD & CI & ARI & Silh. & DB \\
\midrule
Untransformed       & $92.99$ & $7.34$ & $[71.76,\,96.86]$ & $\mathbf{0.864}$ & $0.709$ & $0.528$ \\
Min--max            & $93.42$ & $3.47$ & $[93.91,\,93.91]$ & $0.847$ & $0.711$ & $0.503$ \\
$z$-score           & $92.48$ & $4.03$ & $[88.14,\,93.91]$ & $0.823$ & $0.693$ & $0.542$ \\
Rank transform      & $93.64$ & $2.73$ & $[93.91,\,93.91]$ & $0.849$ & $0.671$ & $0.601$ \\
QuantileTransformer & $93.91$ & $\mathbf{0.00}$ & $[93.91,\,93.91]$ & $0.852$ & $0.667$ & $0.610$ \\
\botrule
\end{tabular}
\end{table}

\begin{table}[h]
\centering
\caption{Protocol B: stratified five-fold cross-validation, 20 seeds, 100 runs.
Accuracy on held-out folds, Rose-4400.}\label{tab:protocolB}
\begin{tabular}{@{}lccc@{}}
\toprule
Representation & Accuracy & SD & CI \\
\midrule
Untransformed       & $\mathbf{93.65}$ & $6.27$ & $[70.96,\,97.50]$ \\
QuantileTransformer & $93.40$ & $3.67$ & $[93.30,\,94.43]$ \\
Rank transform      & $92.82$ & $5.40$ & $[69.65,\,94.43]$ \\
Min--max            & $92.49$ & $5.71$ & $[70.96,\,94.43]$ \\
$z$-score           & $91.61$ & $4.82$ & $[72.53,\,94.43]$ \\
\botrule
\end{tabular}
\end{table}

\begin{table}[h]
\centering
\caption{Paired $t$-tests across 100 seeds, Protocol A, Rose-4400. A positive
mean difference favours representation A.}\label{tab:ttest}
\begin{tabular}{@{}llcccc@{}}
\toprule
Representation A & Representation B & Diff.\ (pp) & $t$ & $p$ & $d_z$ \\
\midrule
Rank transform      & Untransformed & $+0.65$ & $0.821$ & $0.413$ & $0.08$ \\
Min--max            & Untransformed & $+0.43$ & $0.568$ & $0.571$ & $0.06$ \\
QuantileTransformer & Untransformed & $+0.92$ & $1.251$ & $0.214$ & $0.13$ \\
Untransformed       & $z$-score     & $+0.51$ & $0.730$ & $0.467$ & $0.07$ \\
Rank transform      & $z$-score     & $+1.15$ & $2.336$ & $0.022$ & $0.23$ \\
QuantileTransformer & $z$-score     & $+1.43$ & $3.544$ & $0.001$ & $0.35$ \\
\botrule
\end{tabular}
\end{table}

Three conclusions follow. First, no transformation improves mean accuracy over
the untransformed descriptor at any conventional significance level; the
largest such comparison gives $p=0.214$ with $d_z=0.13$. Second, under the
leakage-free Protocol B the untransformed descriptor is in fact the best
performer, the opposite of what earlier reports on this pipeline claim. Third,
and this is the one respect in which the transformations earn their place, they
reduce sensitivity to initialisation sharply: the standard deviation across
seeds falls from $7.34$ percentage points untransformed to $2.73$ under the
rank transform and to $0.00$ under \texttt{QuantileTransformer}, which returned
the same partition at all 100 seeds. Practitioners who need reproducible output
from a single run should apply the transform; those who report distributions
across seeds gain nothing from it.

\subsection{Comparison of clustering algorithms}\label{sec:algos}

Table~\ref{tab:algorithms} compares four clustering algorithms on the same
descriptor, each tuned by silhouette under identical conditions.

\begin{table}[h]
\centering
\caption{Clustering algorithms on the seven-dimensional descriptor, Rose-4400.
Hyperparameters selected by silhouette. Deterministic methods have zero
standard deviation by construction. Runtime is per run on a laptop
CPU.}\label{tab:algorithms}
\begin{tabular}{@{}llccc@{}}
\toprule
Algorithm & Selected configuration & Accuracy & SD & Time (s) \\
\midrule
Agglomerative    & ward linkage & $\mathbf{95.84}$ & $0.00$ & $0.20$ \\
$K$-means        & $k$-means++, $K=4$ & $93.33$ & $6.23$ & $0.25$ \\
Gaussian mixture & full covariance, $k$-means init & $93.33$ & $6.23$ & $0.26$ \\
Gaussian mixture & spherical, random init & $75.71$ & $12.71$ & $0.24$ \\
DBSCAN           & $\epsilon=0.310$, min.\ samples $=50$ & $72.77$ & $0.00$ & $0.09$ \\
\botrule
\end{tabular}
\end{table}

Agglomerative clustering with ward linkage outperforms $K$-means by $2.5$
percentage points and is deterministic, at comparable cost. We report this
because it bears on how the contribution should be read: the strength of the
approach lies in the descriptor, not in any particular clustering algorithm,
and a practitioner should prefer ward linkage here.

Two implementation details are worth recording, since either would otherwise
produce misleading numbers. The default initialisation of the Gaussian mixture
implementation in scikit-learn is $k$-means with the same random seed, so the
mixture converges to the same partition as $K$-means and that comparison
carries no information; the row with random initialisation is the informative
one. Second, an unconstrained silhouette search for DBSCAN selects a very small
$\epsilon$ that splits the data into more than twenty clusters of duplicated
points, achieving a degenerate silhouette of $1.0$. The descriptor contains
many tied values, so this outcome is not incidental. We therefore restricted
the DBSCAN grid to configurations yielding exactly $K$ clusters.

\subsection{Verification of photometric invariance}\label{sec:invariance-exp}

We applied gamma correction $\varphi(a)=255\,(a/255)^{\gamma}$ with
$\gamma\in\{0.7,1.3\}$ to every pixel of all 4400 images and recomputed
descriptors on the transformed grid $\{\varphi(100),\varphi(150),
\varphi(200)\}$. The maximum absolute deviation from the original descriptors,
over all images and coordinates, was $0$ to machine precision, confirming
Proposition~\ref{prop:invariance}. Holding the sampling grid fixed at
$(100,150,200)$ instead, clustering accuracy was unchanged while the mean
coordinatewise displacement was $0.097$ for $\gamma=0.7$.

\subsection{Chebyshev bounds and observed confusions}\label{sec:bounds-exp}

Table~\ref{tab:cheb-bounds} evaluates \eqref{eq:cheb} for every pair of rose
classes, minimised over the seven coordinates, and
Table~\ref{tab:confusion-rose} gives the confusion matrix at the best seed. The
pair with the loosest bound, pink versus red at $15.4\%$, accounts for almost
all misassignments; the pair with the tightest bound, red versus white at
$1.4\%$, is never confused. The ordering induced by the bound agrees with the
empirical error structure.

\begin{table}[h]
\centering
\caption{Chebyshev bounds \eqref{eq:cheb} for Rose-4400, minimised over the
seven coordinates.}\label{tab:cheb-bounds}
\begin{tabular}{@{}llccc@{}}
\toprule
Class pair & Coordinate & $\Delta_k$ & $\sigma_k$ & $\varepsilon_k$ \\
\midrule
red vs.\ white    & $G_{200}$ & $0.249$ & $0.015$ & $1.4\%$ \\
white vs.\ yellow & $B_{150}$ & $0.391$ & $0.042$ & $4.6\%$ \\
red vs.\ yellow   & $G_{150}$ & $0.600$ & $0.072$ & $5.7\%$ \\
pink vs.\ yellow  & $B_{150}$ & $0.497$ & $0.063$ & $6.5\%$ \\
pink vs.\ white   & $R_{200}$ & $0.141$ & $0.025$ & $12.3\%$ \\
pink vs.\ red     & $B_{150}$ & $0.447$ & $0.088$ & $15.4\%$ \\
\botrule
\end{tabular}
\end{table}

\begin{table}[h]
\centering
\caption{Confusion matrix on Rose-4400, $k$-means++ at the best seed, after
Hungarian matching. Rows are true classes.}\label{tab:confusion-rose}
\begin{tabular}{@{}lcccc@{}}
\toprule
 & pink & red & white & yellow \\
\midrule
pink   & $1015$ & $85$  & $0$    & $0$ \\
red    & $138$  & $962$ & $0$    & $0$ \\
white  & $0$    & $0$   & $1100$ & $0$ \\
yellow & $0$    & $0$   & $45$   & $1055$ \\
\botrule
\end{tabular}
\end{table}

\subsection{Ablations}\label{sec:ablation}

%% PENTING: tabel di bawah dihitung di lingkungan berbeda dari tabel lain,
%% sehingga baris pertama membaca 92.76 sedangkan Tabel 3 membaca 92.99.
%% Jalankan ulang di mesin Anda dan ganti seluruh isi tabel ini.
%% Sel kode untuk baseline_mandiri.ipynb:
%%
%%   sets = {
%%     "(100,150,200) tujuh koordinat": COLS7,
%%     "(100,150,200) grid penuh": [(c,t) for c in "RGB" for t in (100,150,200)],
%%     "(50,100,150)":  [(c,t) for c in "RGB" for t in (50,100,150)],
%%     "(150,200,250)": [(c,t) for c in "RGB" for t in (150,200,250)],
%%     "(100,200)":     [(c,t) for c in "RGB" for t in (100,200)],
%%     "(64,128,192)":  [(c,t) for c in "RGB" for t in (64,128,192)],
%%   }
%%   flat = imgs.reshape(len(imgs), -1, 3).astype(np.int16)
%%   ch = {"R": flat[:,:,0], "G": flat[:,:,1], "B": flat[:,:,2]}
%%   rows = []
%%   for name, pairs in sets.items():
%%       X = np.column_stack([(ch[c] <= t).mean(axis=1) for c, t in pairs])
%%       r, _ = protocol_A(X, y, labs, K, n_seeds=N_SEEDS, tag=name)
%%       r["d"] = X.shape[1]; rows.append(r)
%%   pd.DataFrame(rows).to_csv(f"{OUTDIR}/tabel_ablasi_level.csv", index=False)

Table~\ref{tab:ablation-levels} varies the sampling grid. The choice of levels
matters considerably: shifting the grid downwards to $(50,100,150)$ costs more
than 25 percentage points, and an adaptive grid at the global $25$th, $50$th
and $75$th percentiles performs no better. The grid $(150,200,250)$ is
competitive with our choice and attains the highest single-seed accuracy. We
therefore present $(100,150,200)$ as an empirically adequate choice for this
dataset, not as a recommendation that transfers.

\begin{table}[h]
\centering
\caption{Effect of the sampling grid, Rose-4400, Protocol A, 100
seeds.}\label{tab:ablation-levels}
\begin{tabular}{@{}lcccc@{}}
\toprule
Sampling grid & $d$ & Accuracy & SD & Best seed \\
\midrule
$(100,150,200)$, seven coordinates  & $7$ & $92.76$ & $7.92$  & $96.86$ \\
$(150,200,250)$, full grid          & $9$ & $92.96$ & $8.93$  & $\mathbf{98.07}$ \\
$(100,150,200)$, full grid          & $9$ & $89.22$ & $12.19$ & $96.86$ \\
$(100,200)$, full grid              & $6$ & $86.23$ & $14.61$ & $96.86$ \\
Adaptive percentiles $(59,141,196)$ & $9$ & $73.36$ & $11.60$ & $96.86$ \\
$(64,128,192)$, full grid           & $9$ & $71.04$ & $10.80$ & $96.86$ \\
$(50,100,150)$, full grid           & $9$ & $67.59$ & $6.21$  & $96.86$ \\
\botrule
\end{tabular}
\end{table}

%% PENTING: tabel di bawah dihitung dengan resize bikubik (default PIL),
%% sedangkan seluruh tabel lain memakai resize bilinear (default cv2), sehingga
%% baris p=32 membaca 90.78 dan bukan 92.99 seperti Tabel 3. Jalankan ulang
%% dengan cv2.resize lalu ganti isi tabel ini sebelum submit. Kesimpulan
%% paragrafnya tidak berubah, karena seluruh selisih antar-resolusi berada di
%% bawah satu simpangan baku.

Table~\ref{tab:ablation-size} varies the resize resolution over a factor of
eight. Accuracy is flat across the range: the spread between the best and worst
resolution is $2.7$ percentage points against standard deviations between
$7.2$ and $9.2$, so no resolution differs from any other by more than a third
of a standard deviation. The adjusted Rand index and the silhouette coefficient
are likewise stable, varying by $0.046$ and $0.012$ respectively.

This is the behaviour Definition~\ref{def:measure} predicts. The descriptor
depends on an image only through its colour distribution, and bilinear
downsampling is a local averaging operation that leaves that distribution
substantially intact; it removes spatial detail, which the descriptor never
used. We adopt $p=32$ because it is the smallest resolution at which no
degradation is observable, and because the cost of the extraction stage scales
as $O(Np^2)$ by Proposition~\ref{prop:complexity} while the descriptor
dimension stays at seven regardless of $p$. A practitioner working at
$p=128$ pays sixteen times the extraction cost for no measurable gain.

\begin{table}[h]
\centering
\caption{Effect of resize resolution, Rose-4400, Protocol A, 100
seeds.}\label{tab:ablation-size}
\begin{tabular}{@{}lcccc@{}}
\toprule
Resolution & Accuracy & SD & ARI & Silh. \\
\midrule
$16\times16$   & $91.30$ & $7.23$ & $0.823$ & $\mathbf{0.715}$ \\
$32\times32$   & $90.78$ & $7.77$ & $0.821$ & $0.705$ \\
$64\times64$   & $89.42$ & $9.18$ & $0.812$ & $0.703$ \\
$128\times128$ & $\mathbf{92.11}$ & $8.63$ & $\mathbf{0.857}$ & $0.707$ \\
\botrule
\end{tabular}
\end{table}

\subsection{Flowers102-C}\label{sec:flowers}

Table~\ref{tab:flowers-descriptors} repeats the descriptor comparison on
Flowers102-C under Protocol A. Every method performs far worse than on
Rose-4400, and the ordering changes: the RGB histogram now edges ahead of the
ECDF descriptor, and the differences among all four descriptors lie within
roughly one standard deviation of each other. No colour representation
succeeds here, which is the point of including the dataset.

\begin{table}[h]
\centering
\caption{Protocol A on Flowers102-C, 100 seeds, $K=5$.}
\label{tab:flowers-descriptors}
\begin{tabular}{@{}lcccccc@{}}
\toprule
Descriptor & $d$ & Accuracy & SD & ARI & NMI & Silh. \\
\midrule
RGB histogram, $4^3$ bins & $64$ & $45.23$ & $6.78$ & $0.147$ & $0.237$ & $0.058$ \\
ECDF, full grid           & $9$  & $43.85$ & $2.93$ & $0.129$ & $0.204$ & $0.209$ \\
Colour moments            & $9$  & $42.60$ & $4.76$ & $0.135$ & $0.196$ & $0.208$ \\
ECDF, seven coordinates   & $7$  & $42.22$ & $4.19$ & $0.128$ & $0.205$ & $0.220$ \\
\botrule
\end{tabular}
\end{table}

The supervised ceiling confirms that the limitation is informational rather
than algorithmic. On the seven-dimensional descriptor, a linear SVM reaches
$56.97\%$ and linear discriminant analysis $59.16\%$ under five-fold
cross-validation, against $100\%$ on Rose-4400. Even with full access to the
labels, colour cannot separate these five categories.

Table~\ref{tab:flowers-confusion} gives the confusion matrix at the best seed.
The brown category, which has 21 images against 40 for the others, is never
recovered: not one of its images is assigned to its own cluster. Yellow acts as
an attractor for every class.

\begin{table}[h]
\centering
\caption{Confusion matrix on Flowers102-C, $k$-means++ at the best seed, after
Hungarian matching. Rows are true classes.}\label{tab:flowers-confusion}
\begin{tabular}{@{}lccccc@{}}
\toprule
 & blue & brown & purple & red & yellow \\
\midrule
blue   & $16$ & $4$ & $8$  & $0$  & $12$ \\
brown  & $6$  & $0$ & $2$  & $3$  & $10$ \\
purple & $1$  & $3$ & $24$ & $3$  & $9$ \\
red    & $2$  & $2$ & $5$  & $20$ & $11$ \\
yellow & $3$  & $1$ & $3$  & $0$  & $33$ \\
\botrule
\end{tabular}
\end{table}

Because the classes are unbalanced, accuracy after Hungarian matching is biased
towards the larger categories; we therefore report ARI and NMI alongside it,
neither of which depends on the matching. With 181 images, five-fold
cross-validation leaves roughly 36 images per held-out fold, so the standard
deviations under Protocol B are correspondingly large: $44.68\pm5.15$
untransformed, $49.16\pm4.86$ with min--max scaling, $46.65\pm5.47$ with the
rank transform.

Two further observations concern the clustering algorithms. Agglomerative
clustering, which was the strongest method on Rose-4400, is the weakest here at
$24.86\%$: silhouette selects average linkage, which produces one dominant
cluster and four near-empty ones with a deceptively high silhouette of $0.334$.
Internal validity indices reward that degenerate solution precisely because the
classes are not separated in the feature space. DBSCAN returned no
configuration yielding exactly five clusters anywhere in the grid. Both
outcomes are symptoms of the same underlying condition, quantified next.

%%%%%%%%%%%%%%%%%%%%%%%%%%%%%%%%%%%%%%%%%%%%%%%%%%%%%%%%%%%%%%%%%%%%%%%%
\section{Discussion}\label{sec:discussion}

\subsection{When does a colour-only descriptor work?}\label{sec:failure}

The gap between our two datasets is a property of the data, not of the
algorithm. Let $B/W$ denote the ratio of between-class to within-class variance
in mean-RGB space. On Rose-4400 this ratio is $7.481$; on Flowers102-C it is
$0.361$, roughly twenty times smaller. Table~\ref{tab:flowers-stats} shows why:
the mean red-channel values of the purple and red categories differ by fewer
than five intensity levels while their within-class standard deviations exceed
$27$, and the blue category has a near-neutral mean. Threshold-derived dominant-colour labels
applied to flower photographs with dominant green backgrounds do not
correspond to separated regions of whole-image colour space.

\begin{table}[h]
\centering
\caption{Class statistics in mean-RGB space,
Flowers102-C.}\label{tab:flowers-stats}
\begin{tabular}{@{}lccc@{}}
\toprule
Class & $n$ & Mean RGB & SD RGB \\
\midrule
blue   & $40$ & $(90.7,\ 94.5,\ 106.8)$  & $(30.0,\ 30.6,\ 35.4)$ \\
brown  & $21$ & $(114.7,\ 91.4,\ 82.7)$  & $(25.5,\ 20.3,\ 26.0)$ \\
purple & $40$ & $(139.7,\ 95.0,\ 114.9)$ & $(27.9,\ 23.0,\ 26.0)$ \\
red    & $40$ & $(135.3,\ 79.3,\ 67.5)$  & $(34.0,\ 25.4,\ 29.0)$ \\
yellow & $40$ & $(103.8,\ 103.2,\ 75.7)$ & $(22.1,\ 21.3,\ 23.6)$ \\
\botrule
\end{tabular}
\end{table}

Theorem~\ref{thm:separability}(ii) detects the same condition, and does so
more sharply. Table~\ref{tab:flowers-cheb} evaluates the bound
\eqref{eq:cheb} for all ten class pairs of Flowers102-C, minimised over the
seven coordinates. Every bound is vacuous: for each pair, the within-class
standard deviation exceeds half the separation of the class means along every
available coordinate, so $4\sigma_k^2/\Delta_k^2>1$ throughout and the theorem
guarantees nothing. On Rose-4400 the same computation returned bounds between
$1.4\%$ and $15.4\%$ (Table~\ref{tab:cheb-bounds}). The theory therefore
predicts both the success on one dataset and the failure on the other, from
quantities estimable before any clustering is run.

\begin{table}[h]
\centering
\caption{Chebyshev bounds \eqref{eq:cheb} for Flowers102-C, minimised over the
seven coordinates. Every bound is vacuous.}\label{tab:flowers-cheb}
\begin{tabular}{@{}llccc@{}}
\toprule
Class pair & Coordinate & $\Delta_k$ & $\sigma_k$ & $\varepsilon_k$ \\
\midrule
blue vs.\ purple   & $R_{150}$ & $0.267$ & $0.171$ & $\ge 100\%$ \\
purple vs.\ yellow & $R_{150}$ & $0.250$ & $0.171$ & $\ge 100\%$ \\
red vs.\ yellow    & $R_{200}$ & $0.244$ & $0.147$ & $\ge 100\%$ \\
blue vs.\ red      & $R_{150}$ & $0.236$ & $0.151$ & $\ge 100\%$ \\
brown vs.\ purple  & $B_{150}$ & $0.195$ & $0.166$ & $\ge 100\%$ \\
purple vs.\ red    & $B_{150}$ & $0.191$ & $0.166$ & $\ge 100\%$ \\
blue vs.\ brown    & $B_{150}$ & $0.185$ & $0.189$ & $\ge 100\%$ \\
blue vs.\ yellow   & $B_{150}$ & $0.178$ & $0.189$ & $\ge 100\%$ \\
brown vs.\ red     & $R_{200}$ & $0.164$ & $0.147$ & $\ge 100\%$ \\
brown vs.\ yellow  & $R_{200}$ & $0.079$ & $0.096$ & $\ge 100\%$ \\
\botrule
\end{tabular}
\end{table}

This suggests a practical criterion. Computing $B/W$ costs one pass over the
data and requires no clustering, and the bounds $\varepsilon_k$ cost little
more. Where $B/W$ is well above unity and the bounds are informative, a
low-dimensional colour descriptor is likely to suffice; where $B/W$ falls below
unity and every bound is vacuous, texture or shape information is necessary and
no amount of tuning of a colour-only method will compensate. The two
diagnostics agree on both of our datasets, and neither requires labels beyond
the class assignment used to estimate the moments.

A cautionary observation follows from Section~\ref{sec:flowers}. When classes
are not separated in the feature space, internal validity indices do not merely
fail to help, they actively mislead: silhouette selected average linkage for
agglomerative clustering on Flowers102-C, awarding $0.334$ to a degenerate
partition that scores $24.86\%$ accuracy, while $K$-means with a lower
silhouette of $0.222$ scored $42.01\%$. Hyperparameter selection by internal
index, which is the right choice for a fair label-free comparison, is only
trustworthy in the regime the criterion above identifies.

\subsection{Clustering versus classification}\label{sec:gap}

The descriptor supports perfect linear classification
(Table~\ref{tab:supervised}), yet the best clustering result is $95.84\%$ and
$K$-means reaches only $93.33\%$. The gap is attributable to the clustering
objectives rather than to the representation. $K$-means seeks spherical
clusters of comparable size and cannot exploit the anisotropic, elongated
geometry that the linear classifiers use; ward linkage, which merges by
variance increase rather than by distance to a centroid, recovers part of the
difference. Spectral clustering on a suitable affinity is a natural next step.

\subsection{Limitations}\label{sec:limitations}

The descriptor uses colour only and fails on texture-dependent tasks, as
Flowers102-C demonstrates: no colour descriptor we tested exceeded $45.23\%$
there, and the supervised ceiling was $59.16\%$. The intensity levels are fixed and were chosen
empirically on one dataset; Table~\ref{tab:ablation-levels} shows the choice
matters, and we do not claim it transfers. The number of clusters must be
specified in advance. Preprocessing resolution, by contrast, is not a
sensitivity: Table~\ref{tab:ablation-size} shows accuracy flat from
$16\times16$ to $128\times128$. The invariance of Proposition~\ref{prop:invariance} holds
for a photometric transformation applied uniformly across images; per-image
illumination variation is not covered and remains a source of within-class
variance. Finally, our deep-feature baseline uses a single frozen
ImageNet-pretrained backbone; that comparison speaks to inductive bias, not to
what a network trained for this task could achieve.

%%%%%%%%%%%%%%%%%%%%%%%%%%%%%%%%%%%%%%%%%%%%%%%%%%%%%%%%%%%%%%%%%%%%%%%%
\section{Conclusion}\label{sec:conclusion}

Seven numbers per image, obtained by evaluating per-channel empirical
distribution functions at three intensity levels, separate four rose colour
classes perfectly under linear classification and support clustering at
$95.84\%$ with ward linkage and $92.99\pm7.34\%$ with $K$-means over 100
initialisations. The descriptor outperforms colour moments, a 64-dimensional
RGB histogram, and a frozen ImageNet-pretrained ResNet-50, the last by a margin
we attribute to a mismatch of inductive bias rather than to any general
superiority over learned representations.

We proved and verified exact invariance under monotone photometric
transformations, and replaced an incorrect separability claim common in this
literature with an exact result under a margin condition together with a
distribution-free bound whose ordering matches the observed confusion
structure. Two negative results accompany the positive ones: a rank transform
of the descriptor leaves mean accuracy unchanged while sharply reducing
sensitivity to initialisation, and $K$-means is not the best clustering
algorithm available for this representation. On a five-category colour subset of Flowers102
every colour descriptor fails, and the same Chebyshev bound that predicts the
error structure on Rose-4400 is vacuous for all ten class pairs there. Together
with the between-class to within-class colour variance ratio, which falls from
$7.481$ to $0.361$ between the two datasets, this gives a diagnostic that
distinguishes the regime in which the approach succeeds from the one in which
it does not, computable before committing to a colour-only method.

%%%%%%%%%%%%%%%%%%%%%%%%%%%%%%%%%%%%%%%%%%%%%%%%%%%%%%%%%%%%%%%%%%%%%%%%
\backmatter

\bmhead{Acknowledgements}
This research was supported by the Department of Mathematics and the School of
Electrical Engineering and Informatics, Institut Teknologi Bandung, and by the
Department of Mathematics Education, Universitas Pendidikan Indonesia.

\section*{Declarations}

\begin{itemize}
\item \textbf{Funding.} \TODO{nomor hibah, atau nyatakan tidak ada pendanaan
khusus}
\item \textbf{Competing interests.} The authors declare no competing interests.
\item \textbf{Ethics approval.} Not applicable.
\item \textbf{Consent to participate.} Not applicable.
\item \textbf{Consent for publication.} All authors consent to publication.
\item \textbf{Data availability.} The Flowers102 dataset is publicly available
\cite{ref:nilsback2008}, and the rose collection is publicly available on
Kaggle \cite{ref:rosekaggle}. A single notebook reproduces every table in
this paper end to end, from downloading both collections through extraction to
the reported metrics; it is available from the corresponding author on
reasonable request. \TODO{lebih baik: unggah notebook ke repositori publik dan
ganti kalimat terakhir dengan tautan permanennya, misalnya DOI Zenodo}
\item \textbf{Code availability.} \TODO{unggah kode ke repositori publik dan
cantumkan tautan permanen di sini}
\item \textbf{Author contribution.} \TODO{isi sesuai kontribusi masing-masing
penulis}
\end{itemize}

\begin{appendices}

\section{Illustrative example on a six-image dataset}\label{secA1}

\TODO{Opsional, dan kami menyarankan menghapus bagian ini seluruhnya. Enam
citra tidak mendukung kesimpulan apa pun, reviewer sebelumnya sudah memintanya
dipersingkat atau dipindahkan, dan berkas ekstraksinya memakai $p=3$ yaitu
sembilan piksel per citra, bukan $32\times32$ seperti dinyatakan di naskah
lama.}

\end{appendices}

\bibliography{sn-bibliography}

\end{document}
